\section{Label observables under restriction to two groups of copies}
\label{sec:grouped_label_entropy}

The subgroup dimension inequalities of \cite[Lemma~6.1(4)]{OpenAI2026AreaLaw}
give matrix inequalities for the logarithmic label observables. These are
used in Section~7 to compare whole-copy and subgroup labels before measuring
an excitation component.

\begin{theorem}[Grouped-copy label bounds]
\label{thm:grouped_label_entropy_bounds}
\lean{TensorPower.groupedCopies_labelEntropy_bounds}
\leanok
Let $m+r=k$, fix a partition of the $k$ copies into specified groups of
sizes $m$ and $r$, and let $\phi:S_k\to\operatorname{Perm}(X)$ be any
permutation representation on a finite coordinate set $X$. For the
restrictions to the two groups and for the whole group, let
\[
 F_g=\sum_\alpha(\log d_\alpha)P^g_\alpha,\qquad
 F_b=\sum_\beta(\log d_\beta)P^b_\beta,\qquad
 F_w=\sum_\lambda(\log d_\lambda)P^w_\lambda
\]
be the central label observables on $\mathbb C^X$. Then, on the whole space,
\begin{align}
 F_g+F_b\le F_w\le F_g+F_b+\log\binom{k}{r}\,I.
 \label{eq:grouped_label_entropy_bounds}
\end{align}
The statement includes $m=0$ and $r=0$.
\end{theorem}
\begin{proof}
\leanok
The three families of central label projections commute. Their products
$P^g_\alpha P^b_\beta P^w_\lambda$ form an orthogonal resolution of the
identity. On every nonzero product, the two subgroup dimension bounds give
\[
 d_\alpha d_\beta\le d_\lambda
 \le\binom{k}{r}d_\alpha d_\beta.
\]
All irreducible dimensions are positive, and $r\le k$ implies
$\binom{k}{r}>0$. Taking logarithms proves both inequalities on this
joint component. Vanishing products make no contribution. Summing over the
joint resolution proves \eqref{eq:grouped_label_entropy_bounds}.
\end{proof}

The ensuing good-auxiliary estimate also uses a dimension bound for labels
which actually occur on the bad copies and preservation of the whole-copy
auxiliary labels by a physical excitation projection. Those are separate
assertions.
